\documentclass[11pt]{article}
\usepackage{amsmath,amssymb,amsthm}
\usepackage{hyperref}

\newtheorem{theorem}{Theorem}
\newtheorem{lemma}{Lemma}

\newcommand{\R}{\mathbb R}
\newcommand{\C}{\mathbb C}
\newcommand{\sech}{\operatorname{sech}}
\newcommand{\sgn}{\operatorname{sgn}}
\newcommand{\supp}{\operatorname{supp}}
\newcommand{\calF}{\mathcal F}

\title{Endpoint Failure for an Unbounded-Group Transference Principle}
\author{Run \texttt{fa\_banach\_001}, agent \texttt{agent\_lane\_08}}
\date{2026-07-19}

\begin{document}
\maketitle

\section*{Source and Target}

The source paper is Markus Haase, \emph{A Transference Principle for General
Groups and Functional Calculus on UMD Spaces}, arXiv:0807.3906. In Section 3,
Theorem 3.2 proves a transference estimate for a \(C_0\)-group \(U\) satisfying
\[
  \|U(s)\| \le M\cosh(\omega_0 s)
\]
when the measure weight parameter \(\omega\) is strictly larger than
\(\omega_0\). In Remark 3.3 the source says that this strict inequality is
essential in the proof and that the endpoint \(\omega=\omega_0\) is not
expected.

This packet proves that the endpoint estimate indeed fails, already in the
scalar Hilbert-space case.

\section*{Endpoint Form Refuted}

At the endpoint, the analogue of Haase's Theorem 3.2 would assert the following
kind of estimate: for \(p=2\), \(X=\C\), a \(C_0\)-group \(U\) with
\(\|U(s)\|\le M\cosh(\omega_0s)\), and all \(\mu\in M_{\omega_0}(\R)\),
\[
  \left|\int_\R U(s)\,\mu(ds)\right|
  \le C M^2\, \|L_{\mu_{\omega_0}}\|_{L^2(\R)\to L^2(\R)},
  \qquad
  \mu_{\omega_0}(ds)=\cosh(\omega_0s)\mu(ds).
\]
Here \(L_\nu\) denotes convolution by the finite measure \(\nu\).

\section*{Proof Intuition}

Take the scalar group \(U(s)=e^{\alpha s}\), whose growth type is exactly
\(\alpha>0\). If \(\nu=\mu_\alpha\), then the transferred scalar operator is
pairing \(\nu\) with
\[
  \frac{e^{\alpha s}}{\cosh(\alpha s)}=1+\tanh(\alpha s).
\]
For scalar \(L^2\), the convolution norm of \(\nu\) is just the supremum norm
of its Fourier transform. Thus the endpoint estimate would say that pairing
with \(\tanh(\alpha s)\) is controlled by \(\|\widehat\nu\|_\infty\), at least
for measures of total mass zero.

This is false. The Fourier transform of \(\tanh(\alpha s)\) has a principal
value \(1/\xi\) singularity at frequency \(0\). Smooth Fourier multipliers
bounded by one and supported on \(1/N\lesssim |\xi|\lesssim 1\) therefore pair
with it like \(\int_{1/N}^1 d\xi/\xi\), which grows as \(\log N\).

\section*{Fourier Convention}

We use
\[
  \widehat f(\xi)=\int_\R e^{-is\xi}f(s)\,ds,\qquad
  f(s)=\frac1{2\pi}\int_\R e^{is\xi}\widehat f(\xi)\,d\xi.
\]

\begin{lemma}\label{lem:tanh}
For \(\alpha>0\), in the sense of tempered principal-value distributions,
\[
  \widehat{\tanh(\alpha\,\cdot)}(\xi)
  =
  \frac{\pi}{i\alpha\sinh(\pi\xi/(2\alpha))}.
\]
\end{lemma}

\begin{proof}
The derivative of \(\tanh(\alpha s)\) is \(\alpha\sech^2(\alpha s)\).
The standard contour computation gives
\[
  \calF[\alpha\sech^2(\alpha\,\cdot)](\xi)
  =
  \frac{\pi\xi/\alpha}{\sinh(\pi\xi/(2\alpha))}.
\]
Since differentiation corresponds to multiplication by \(i\xi\), division by
\(i\xi\) gives the displayed formula away from zero. The distribution is odd
and purely imaginary, so there is no additional delta mass at zero.
\end{proof}

\begin{theorem}
Let \(\alpha>0\). The endpoint transference estimate with
\(\omega=\omega_0=\alpha\) fails for the scalar Hilbert space \(X=\C\), \(p=2\),
and \(U(s)=e^{\alpha s}\).
\end{theorem}

\begin{proof}
The group \(U(s)=e^{\alpha s}\) satisfies
\[
  |U(s)|\le 2\cosh(\alpha s),\qquad s\in\R,
\]
and its growth type is exactly \(\alpha\).

Assume, toward a contradiction, that an endpoint estimate holds with some
finite constant \(C\) for this fixed group. Put
\[
  a=\frac{\pi}{2\alpha}.
\]
Choose \(\delta>0\) so small that
\[
  |\sinh(a\xi)|\le 2a|\xi|
  \qquad (|\xi|\le \delta).
\]
For each sufficiently large \(N\), let \(\chi_N\in C_c^\infty(\R)\) be even,
\(0\le\chi_N\le1\), with
\[
  \chi_N(\xi)=1
  \quad\text{for}\quad
  2/N\le |\xi|\le \delta/2,
  \qquad
  \supp \chi_N\subset\{1/N\le |\xi|\le \delta\}.
\]
Define
\[
  m_N(\xi)=-i\,\sgn(\xi)\chi_N(\xi).
\]
Then \(m_N\in C_c^\infty(\R)\), \(m_N(0)=0\), and
\(\|m_N\|_\infty\le1\). Let
\[
  h_N(s)=\frac1{2\pi}\int_\R e^{is\xi}m_N(\xi)\,d\xi,
  \qquad
  \nu_N(ds)=h_N(s)\,ds .
\]
Since \(m_N\) is smooth and compactly supported, \(h_N\) is a Schwartz
function, hence \(\nu_N\) is a finite measure. Its Fourier transform is
\(\widehat{\nu_N}=m_N\), and its total mass is zero because \(m_N(0)=0\).
Therefore convolution by \(\nu_N\) on scalar \(L^2(\R)\) has norm
\[
  \|L_{\nu_N}\|_{2\to2}=\|m_N\|_\infty\le1.
\]

Now set
\[
  \mu_N(ds)=\sech(\alpha s)\,\nu_N(ds).
\]
Then \(\mu_N\in M_\alpha(\R)\), since
\[
  \int_\R e^{\alpha|s|}\,|\mu_N|(ds)
  \le 2\|\nu_N\|<\infty,
\]
and \((\mu_N)_\alpha=\nu_N\). The transferred scalar operator is
\[
  \int_\R e^{\alpha s}\,\mu_N(ds)
  =
  \int_\R (1+\tanh(\alpha s))\,\nu_N(ds)
  =
  \int_\R \tanh(\alpha s)\,h_N(s)\,ds,
\]
because \(\nu_N(\R)=0\).

By Lemma \ref{lem:tanh} and the definition of \(m_N\),
\[
\begin{aligned}
  \int_\R \tanh(\alpha s)h_N(s)\,ds
  &=
  \frac1{2\pi}
  \left\langle
  \frac{\pi}{i\alpha\sinh(a\xi)},\,m_N(-\xi)
  \right\rangle  \\
  &=
  \frac1{2\alpha}
  \int_\R
  \frac{\chi_N(\xi)}{|\sinh(a\xi)|}\,d\xi .
\end{aligned}
\]
For \(N\) large enough that \(2/N<\delta/2\), the last integral is bounded
below by
\[
  \frac1{2\alpha}
  \int_{2/N\le|\xi|\le\delta/2}
  \frac{d\xi}{2a|\xi|}
  =
  \frac1{2\alpha a}\log\!\left(\frac{\delta N}{4}\right),
\]
which tends to infinity with \(N\).

The hypothesized endpoint transference estimate would instead give
\[
  \left|\int_\R e^{\alpha s}\,\mu_N(ds)\right|
  \le C\,2^2\,\|L_{\nu_N}\|_{2\to2}
  \le 4C
\]
for all \(N\), a contradiction. Hence no endpoint estimate can hold.
\end{proof}

\section*{Verification Notes}

The construction is scalar and uses only Plancherel's theorem for the
\(L^2\) convolution norm. The measures \(\nu_N\) are absolutely continuous
finite measures because their Fourier transforms \(m_N\) are smooth and
compactly supported. The logarithmic divergence comes solely from the
principal-value \(1/\xi\) singularity in the Fourier transform of
\(\tanh(\alpha s)\).

\section*{Novelty and Search Bounds}

The run indexes were searched for arXiv:0807.3906, the source title,
`endpoint transference', `omega\_0', and close variants; no duplicate packet
was found. A bounded web search on 2026-07-19 found later uses and
interpolation variants of Haase's unbounded-group transference principle, but
no exact endpoint counterexample.

\section*{References}

\begin{thebibliography}{9}

\bibitem{Haase0807}
M. Haase,
\emph{A Transference Principle for General Groups and Functional Calculus on
UMD Spaces},
arXiv:0807.3906.

\bibitem{SteinShakarchi}
E. M. Stein and R. Shakarchi,
\emph{Fourier Analysis: An Introduction},
Princeton Lectures in Analysis I, Princeton University Press, 2003.

\end{thebibliography}

\end{document}

